% ============================ METHOD ============================
At every control step \SAS{} reads the body attitude and a motion cue and emits one
velocity per flipper. It pursues two regulation goals: it stabilizes body roll and pitch,
and it releases the robot when it is \emph{stuck}, which means it is commanded forward but
is not advancing. The two terms are formed independently and combined by a
directional-priority rule, so the anti-stuck term can reinforce the stabilization term but
never reverse it.

\SAS{} runs on the four flippers $i\in\{\mathrm{FL},\mathrm{FR},\mathrm{RL},\mathrm{RR}\}$
and emits a per-flipper angular \emph{velocity} command $v_i$. Its inputs are the
gravity-referenced body roll $r$ and pitch $p$ from
the IMU, the four joint angles $\theta_i$, the commanded body velocity $v_\text{cmd}$,
and a scalar motion cue $m$ derived from LiDAR-ICP odometry over a short window
(Sec.~\ref{sec:method-antistuck}).

The command superimposes a stabilization velocity $v^{\mathrm{s}}_i$ and an anti-stuck
velocity $v^{\mathrm{d}}_i$, combined by a directional-priority rule and clamped:
\begin{equation}
  v_i \;=\; \mathrm{limit}\big(\, v^{\mathrm{s}}_i \;+\; \mathrm{cap}(v^{\mathrm{d}}_i,\,v^{\mathrm{s}}_i)\,\big).
  \label{eq:overview}
\end{equation}

\subsection{Attitude stabilization}\label{sec:method-stab}
The dominant term regulates body attitude to level from the IMU alone. The setpoints
$r_0,p_0$ are level, that is zero, and the deviations are $e_r=r_0-r$ and $e_p=p_0-p$. Each
passes through a dead-band of a few degrees, so attitude noise near level produces no
control action. From these we form proportional-integral commands
\begin{equation}
  u_r = K^r_{\mathrm P}\,e_r + K^r_{\mathrm I}\, I_r, \qquad
  u_p = K^p_{\mathrm P}\,e_p + K^p_{\mathrm I}\, I_p,
  \label{eq:pi}
\end{equation}
where $I_{(\cdot)}\!\leftarrow\!\mathrm{clamp}(I_{(\cdot)}+e_{(\cdot)}\Delta t)$ are
anti-windup-limited integrals. The scalar commands map to the four flippers as a front/rear
differential for pitch and a left/right differential for roll:
\begin{equation}
  \begin{bmatrix} v^{\mathrm{s}}_{\mathrm{FL}}\\ v^{\mathrm{s}}_{\mathrm{FR}}\\ v^{\mathrm{s}}_{\mathrm{RL}}\\ v^{\mathrm{s}}_{\mathrm{RR}} \end{bmatrix}
  =
  \begin{bmatrix} -1 & +1\\ -1 & -1\\ +1 & +1\\ +1 & -1 \end{bmatrix}
  \begin{bmatrix} u_p\\ u_r \end{bmatrix}.
  \label{eq:stabmap}
\end{equation}

\subsection{Reactive anti-stuck}\label{sec:method-antistuck}
The subordinate action follows the escape idea of~\cite{azayev2022}: when the robot is
stuck, it drives the flippers to lift the body off the ground and free the robot from the
stuck pose.
It fires whenever the robot is commanded to move yet is not moving:
\begin{equation}
  \text{dig} \;=\; \big(|v_\text{cmd}|>\tau_\text{move}\big)\ \wedge\
  \big(m<\tau_\text{stuck}\big).
  \label{eq:stuck}
\end{equation}
The motion cue $m$ is the standard deviation of the robot's short-horizon displacement
over a sliding window, a \emph{local} measure of whether the body advances, not of where
it is; the window cancels slow drift, so $m$ needs no integrated global pose. It must
measure actual body motion: a track encoder will not do, because a stalled robot can
spin its tracks without advancing. In the current deployment $m$ is the short-window
variance of LiDAR-ICP odometry~\cite{libpointmatcher}. When digging, all
four flippers are driven downward at a lift velocity $\lambda$ so they push against the
terrain and lift the body up, out of the stuck pose; otherwise they retract slowly
toward the flat pose ($-\lambda_\text{low}$):
\begin{equation}
  v^{\mathrm{d}}_i = \begin{cases} +\lambda & \text{dig}\\ -\lambda_\text{low} & \text{otherwise.}\end{cases}
  \label{eq:spread}
\end{equation}
Digging alone can lock up. If the flippers reach their lower position limit without
freeing the robot, the spread the limit permits is zero, the motion cue stays below
$\tau_\text{stuck}$, and the robot would remain stuck indefinitely with its flippers fully
down. A second, longer timer therefore watches the dig itself: after
$T_\text{dig}=\SI{8}{\second}$ of continuous digging without release, the anti-stuck term
is suspended for $T_\text{reset}=\SI{2.5}{\second}$ while the flippers are driven back up
toward the climbing pose, after which digging is permitted again. A dig that cannot free
the robot therefore terminates in a bounded dig--reset cycle rather than a permanent stall.
The reset remains a subordinate demand, so
stabilization still caps it. Stuck detection is armed only after the robot has been
moving for \SI{2}{\second}, so a robot that is merely starting is not mistaken for one that
is stuck.

\subsection{Directional-priority blend and limits}\label{sec:method-blend}
Stabilization takes priority over anti-stuck: the anti-stuck command velocity may reinforce
the stabilization command but must not reverse its direction on any flipper. This is
enforced by the $\mathrm{cap}$ operator, a directional \emph{clamp}: when
$v^{\mathrm{d}}_i$ opposes $v^{\mathrm{s}}_i$ it clamps the anti-stuck command into the band
$[-\alpha|v^{\mathrm{s}}_i|,\,\alpha|v^{\mathrm{s}}_i|]$ set by a fraction $\alpha\in[0,1]$ of the
stabilization demand. It does so only while stabilization is meaningfully active
($|v^{\mathrm{s}}_i|$ above a small dead-band $\beta$), so that attitude noise on a level
robot cannot suppress a dig; otherwise it passes $v^{\mathrm{d}}_i$ through:
\begin{equation}
  \mathrm{cap}(v^{\mathrm{d}}_i,v^{\mathrm{s}}_i)=
  \begin{cases}
    \operatorname{clamp}\big(v^{\mathrm{d}}_i,\ \alpha|v^{\mathrm{s}}_i|\big)
      & \text{if } v^{\mathrm{s}}_i v^{\mathrm{d}}_i<0 \wedge |v^{\mathrm{s}}_i|>\beta,\\[2pt]
    v^{\mathrm{d}}_i & \text{otherwise,}
  \end{cases}
  \label{eq:cap}
\end{equation}
where $\operatorname{clamp}(x,c)$ limits $x$ to $[-c,c]$.

